% Source: https://artofproblemsolving.com/community/c3h2474742_fun_problem
% Stable identity: aops-topic-2474742-q01
Place the cube as $[-6,6]^3$. The square on the face $z=6$ has vertices $(\pm5,\pm5,6)$, and similarly for the other faces; the $24$ vertices are the points with one coordinate $\pm6$ and the other two $\pm5$.

\begin{asy}
usepackage("amsmath");
size(5.5cm);
defaultpen(fontsize(10pt));
triple V=unit((2.0,1.1,0.9));
triple e1=unit(cross((0,0,1),V)); triple e2=cross(V,e1);
pair pr(triple p){return (dot(p,e1),dot(p,e2));}
triple[][] faces;
for(int k=0;k<3;++k) for(int s=-1;s<=1;s+=2){
  triple[] f; real[][] q={{5,5},{-5,5},{-5,-5},{5,-5}};
  for(int i=0;i<4;++i){real[] c=new real[3]; c[k]=6s; c[(k+1)%3]=q[i][0]; c[(k+2)%3]=q[i][1]; f.push((c[0],c[1],c[2]));}
  faces.push(f);}
for(int k=0;k<3;++k) for(int si=-1;si<=1;si+=2) for(int sj=-1;sj<=1;sj+=2){
  int i=k, j=(k+1)%3, m=(k+2)%3;
  triple mk(real a,real bb,real c){real[] r=new real[3]; r[i]=a; r[j]=bb; r[m]=c; return (r[0],r[1],r[2]);}
  faces.push(new triple[]{mk(6si,5sj,5),mk(6si,5sj,-5),mk(5si,6sj,-5),mk(5si,6sj,5)});}
for(int sx=-1;sx<=1;sx+=2) for(int sy=-1;sy<=1;sy+=2) for(int sz=-1;sz<=1;sz+=2)
  faces.push(new triple[]{(6sx,5sy,5sz),(5sx,6sy,5sz),(5sx,5sy,6sz)});
real[] cb={-6,6};
for(int a=0;a<2;++a) for(int c=0;c<2;++c){
  draw(pr((cb[a],cb[c],-6))--pr((cb[a],cb[c],6)),gray(0.7)+dashed);
  draw(pr((cb[a],-6,cb[c]))--pr((cb[a],6,cb[c])),gray(0.7)+dashed);
  draw(pr((-6,cb[a],cb[c]))--pr((6,cb[a],cb[c])),gray(0.7)+dashed);}
for(triple[] f:faces){
  triple g=(0,0,0); for(triple p:f) g+=p;
  if(dot(g,V)>0){
    path pth=pr(f[0]); for(int i=1;i<f.length;++i) pth=pth--pr(f[i]); pth=pth--cycle;
    bool sqf=(f.length==4) && ((f[0].x==f[1].x && f[1].x==f[2].x) || (f[0].y==f[1].y && f[1].y==f[2].y) || (f[0].z==f[1].z && f[1].z==f[2].z));
    pen fc=(f.length==3)?rgb(1,0.85,0.75):(sqf?rgb(0.8,0.88,1):rgb(0.85,0.95,0.85));
    filldraw(pth,fc,black+linewidth(0.7));}}
label("$10$",pr((6,0,-5)),N,fontsize(8pt));
pair qq=pr((5.5,-5,5.5)); draw(qq--qq+(-1.6,1.6),linewidth(0.4)); label("$\sqrt2$",qq+(-1.6,1.6),NW,fontsize(8pt));
\end{asy}

\textbf{The faces.} Every vertex satisfies $|x|,|y|,|z|\le6$, $|x|+|y|\le11$ (and similarly for the other pairs of coordinates), and $|x|+|y|+|z|\le16$. So the polyhedron lies on one side of each of the following planes, and each plane contains at least three non-collinear vertices, so it meets the polyhedron in a face:
\begin{itemize}
\item the $6$ planes such as $z=6$: the drawn $10\times10$ squares, total area $6\cdot100=600$;
\item the $12$ planes such as $y+z=11$ (one for each edge of the cube): the vertices $(\pm5,5,6)$ and $(\pm5,6,5)$ form a rectangle with sides $10$ and $\sqrt{(6-5)^2+(5-6)^2}=\sqrt2$, total area $12\cdot10\sqrt2=120\sqrt2$;
\item the $8$ planes such as $x+y+z=16$ (one for each corner): the vertices $(6,5,5)$, $(5,6,5)$, $(5,5,6)$ form an equilateral triangle with side $\sqrt2$, area $\dfrac{\sqrt3}{4}\cdot2=\dfrac{\sqrt3}{2}$, total $4\sqrt3$.
\end{itemize}
These $26$ polygons fit together edge to edge (each side of a square is a long side of a rectangle, and each short side of a rectangle is a side of a triangle), so they form the whole surface; as a check, $V-E+F=24-48+26=2$. So the surface area is $600+4\sqrt3+120\sqrt2$, while the cube's surface area is $6\cdot12^2=864$. The ratio is
\[\dfrac{600+4\sqrt3+120\sqrt2}{864}=\dfrac{150+\sqrt3+30\sqrt2}{216},\]
and $150,1,30,216$ have no common factor. (Since $1$, $\sqrt2$, $\sqrt3$ are linearly independent over the rationals, this representation is unique.) Hence $a+b+c+d=150+1+30+216=\boxed{397}$.
